En la Tabla~\ref{tab:comparativa} se presenta una
comparación de las principales características
de los trabajos analizados, considerando los
métodos utilizados, el análisis de jugadores
y alineaciones, así como las variables que
se consideran para evaluar su rendimiento.

\begin{table}[htbp]
    \centering
    \caption{Tabla comparativa de los artículos analizados}
    \label{tab:comparativa}
    \small
    \setlength{\tabcolsep}{4pt}
    \renewcommand{\arraystretch}{1.3}

    \begin{tabular}{p{4.3cm}ccccccc}
        \toprule
        \textbf{Trabajo} &
        \textbf{C1} &
        \textbf{C2} &
        \textbf{C3} &
        \textbf{C4} &
        \textbf{C5} &
        \textbf{C6} &
        \textbf{C7} \\
        \midrule

        Chen et al. \cite{Chen17122023}
        & -- & \checkmark & \checkmark
        & \checkmark & -- & \checkmark
        & $\triangle$ \\

        Bonato y Hinds
        \cite{bonato2026networkanalysislinkprediction}
        & $\triangle$ & -- & --
        & \checkmark & $\triangle$ & --
        & $\triangle$ \\

        Josephs y Upton
        \cite{josephs2025hypergraphadjustedplusminus}
        & \checkmark & \checkmark & --
        & \checkmark & \checkmark & $\triangle$
        & -- \\

        Yamada y Fujii
        \cite{yamada2024offensivelineupanalysisbasketball}
        & \checkmark & \checkmark & \checkmark
        & \checkmark & \checkmark & --
        & -- \\

        Devlin y Uminsky
        \cite{devlin2020identifyinggroupcontributionsnba}
        & \checkmark & \checkmark & --
        & -- & \checkmark & $\triangle$
        & -- \\

        \midrule

        \textbf{Propuesta}
        &  &  & 
        &  &  & 
        & \\

        \bottomrule
    \end{tabular}
\end{table}

Las características consideradas para la
comparación son las siguientes:

\begin{itemize}
    \item[C1.] Analiza el rendimiento individual
    de los jugadores.

    \item[C2.] Analiza combinaciones de jugadores
    o alineaciones completas.

    \item[C3.] Utiliza técnicas de clustering
    para clasificar jugadores o alineaciones.

    \item[C4.] Utiliza modelos estadísticos o de
    aprendizaje automático para estimar el
    rendimiento.

    \item[C5.] Considera las interacciones entre
    jugadores de una misma alineación.

    \item[C6.] Considera el rendimiento ofensivo
    y defensivo.

    \item[C7.] Considera las características
    o emparejamientos del equipo contrario
    para evaluar alineaciones.
\end{itemize}

La simbología utilizada en la tabla es:

\begin{itemize}
    \item[$\checkmark$] La característica
    está considerada en el estudio.

    \item[$\triangle$] La característica
    está considerada parcialmente.

    \item[--] La característica no se
    aborda directamente.
\end{itemize}

A partir de la Tabla~\ref{tab:comparativa},
se observa que los trabajos analizados presentan
diferentes enfoques para evaluar el rendimiento
de los jugadores y sus alineaciones. Algunos
estudios se centran en clasificar jugadores
mediante técnicas de agrupamiento, mientras
que otros analizan las contribuciones individuales
y las interacciones entre compañeros mediante
modelos estadísticos y representaciones de grafos.

Sin embargo, los enfoques revisados no integran
de manera conjunta la caracterización de las
jugadoras, el análisis de las combinaciones,
la evaluación ofensiva y defensiva y los
emparejamientos frente al equipo contrario.
